% GENERATED FILE. Edit canonical lessons, then run npm run generate:books.
% canonical-source-hash: fnv1a64:7f1c12e2ef75cf75
% canonical-lessons: TE-C154-itivala, TE-C154-ippatike, TE-C154-gata, TE-C154-okasari, TE-C154-gadupu

\chapter{Recently, Already, Last, Once, To Spend (Time)}
\label{ch:te-recently-already-last-once-to-spend-time}
\hlchaptermodality{154}

\begin{chapteropening}
\textbf{By the end of this chapter:} \emph{I can say what happened, and when: recently, already, last year.}

Complete the last lesson of chapter 154: I can say what happened, and when: recently, already, last year.
\end{chapteropening}

\section[\texorpdfstring{\te{ఇటీవల} (iṭīvala)}{ఇటీవల (iṭīvala)}]{\te{ఇటీవల} (iṭīvala) — recently}

\label{lesson:TE-C154-itivala}



\begin{quote}
Before the new one: say the Telugu for anyone, then the Telugu for anything.
\end{quote}

\subsection*{You'll want to know: \te{ఇటీవల}}
\textbf{\te{ఇటీవల}} — \emph{iṭīvala} — \textquotedblleft{}recently\textquotedblright{}. Telugu's past for \textquotedblleft{}I\textquotedblright{} ends in \textbf{-\te{ఆను}} (\emph{-ānu}): \textbf{\te{నేను} \te{ఇటీవల} \te{ఒక} \te{పుస్తకం} \te{చదివాను}} — \emph{nēnu iṭīvala oka pustakaṁ cadivānu} — \textquotedblleft{}I recently read a book.\textquotedblright{}

\begin{grammarlens}[title={What you've built}]
One more word for this chapter.
\end{grammarlens}

\subsection*{Your turn}
\begin{itemize}
\raggedright
  \item \emph{Say it:} \emph{iṭīvala}
  \item \emph{Say it:} \emph{iṭīvala}, once more
  \item \emph{Say it:} say \emph{iṭīvala}
  \item \emph{Recall:} say the Telugu for anyone, then the Telugu for anything, then say \emph{iṭīvala} again
\end{itemize}

\begin{tcolorbox}[breakable,colback=teal!4,colframe=teal!35!black,title={Before you move on}]
What does \te{ఇటీవల} mean? (\textquotedblleft{}Recently\textquotedblright{}.) Say it once more.
\end{tcolorbox}
\section[\texorpdfstring{\te{ఇప్పటికే} (ippaṭikē)}{ఇప్పటికే (ippaṭikē)}]{\te{ఇప్పటికే} (ippaṭikē) — already}

\label{lesson:TE-C154-ippatike}



\begin{quote}
Before the new one: say the Telugu for anything, then the Telugu for recently.
\end{quote}

\subsection*{You'll want to know: \te{ఇప్పటికే}}
\textbf{\te{ఇప్పటికే}} — \emph{ippaṭikē} — \textquotedblleft{}already\textquotedblright{}. \textbf{\te{నేను} \te{ఇప్పటికే} \te{భోజనం} \te{చేశాను}} — \emph{nēnu ippaṭikē bhōjanaṁ cēśānu} — \textquotedblleft{}I have already eaten.\textquotedblright{}

\begin{grammarlens}[title={What you've built}]
One more word for this chapter.
\end{grammarlens}

\subsection*{Your turn}
\begin{itemize}
\raggedright
  \item \emph{Say it:} \emph{ippaṭikē}
  \item \emph{Say it:} \emph{ippaṭikē}, once more
  \item \emph{Say it:} say \emph{ippaṭikē}
  \item \emph{Recall:} say the Telugu for anything, then the Telugu for recently, then say \emph{ippaṭikē} again
\end{itemize}

\begin{tcolorbox}[breakable,colback=teal!4,colframe=teal!35!black,title={Before you move on}]
What does \te{ఇప్పటికే} mean? (\textquotedblleft{}Already\textquotedblright{}.) Say it once more.
\end{tcolorbox}
\section[\texorpdfstring{\te{గత} (gata)}{గత (gata)}]{\te{గత} (gata) — last, previous}

\label{lesson:TE-C154-gata}



\begin{quote}
Before the new one: say the Telugu for recently, then the Telugu for already.
\end{quote}

\subsection*{You'll want to know: \te{గత}}
\textbf{\te{గత}} — \emph{gata} — \textquotedblleft{}last, previous\textquotedblright{}. \textbf{\te{గత} \te{సంవత్సరం} \te{నేను} \te{హైదరాబాద్} \te{వెళ్ళాను}} — \emph{gata saṁvatsaraṁ nēnu haidarābād veḷḷānu} — \textquotedblleft{}Last year I went to Hyderabad.\textquotedblright{}

\begin{grammarlens}[title={What you've built}]
One more word for this chapter.
\end{grammarlens}

\subsection*{Your turn}
\begin{itemize}
\raggedright
  \item \emph{Say it:} \emph{gata}
  \item \emph{Say it:} \emph{gata}, once more
  \item \emph{Say it:} say \emph{gata}
  \item \emph{Recall:} say the Telugu for recently, then the Telugu for already, then say \emph{gata} again
\end{itemize}

\begin{tcolorbox}[breakable,colback=teal!4,colframe=teal!35!black,title={Before you move on}]
What does \te{గత} mean? (\textquotedblleft{}Last, previous\textquotedblright{}.) Say it once more.
\end{tcolorbox}
\section[\texorpdfstring{\te{ఒకసారి} (okasāri)}{ఒకసారి (okasāri)}]{\te{ఒకసారి} (okasāri) — once, one time}

\label{lesson:TE-C154-okasari}



\begin{quote}
Before the new one: say the Telugu for already, then the Telugu for last.
\end{quote}

\subsection*{You'll want to know: \te{ఒకసారి}}
\textbf{\te{ఒకసారి}} — \emph{okasāri} — \textquotedblleft{}once, one time\textquotedblright{}. \textbf{\te{నేను} \te{ఒకసారి} \te{చెన్నై} \te{వెళ్ళాను}} — \emph{nēnu okasāri cennai veḷḷānu} — \textquotedblleft{}I went to Chennai once.\textquotedblright{}

\begin{grammarlens}[title={What you've built}]
One more word for this chapter.
\end{grammarlens}

\subsection*{Your turn}
\begin{itemize}
\raggedright
  \item \emph{Say it:} \emph{okasāri}
  \item \emph{Say it:} \emph{okasāri}, once more
  \item \emph{Say it:} say \emph{okasāri}
  \item \emph{Recall:} say the Telugu for already, then the Telugu for last, then say \emph{okasāri} again
\end{itemize}

\begin{tcolorbox}[breakable,colback=teal!4,colframe=teal!35!black,title={Before you move on}]
What does \te{ఒకసారి} mean? (\textquotedblleft{}Once, one time\textquotedblright{}.) Say it once more.
\end{tcolorbox}
\section[\texorpdfstring{\te{గడుపు} (gaḍupu)}{గడుపు (gaḍupu)}]{\te{గడుపు} (gaḍupu) — to spend (time)}

\label{lesson:TE-C154-gadupu}



\begin{quote}
Before the new one: say the Telugu for last, then the Telugu for once.
\end{quote}

\subsection*{You'll want to know: \te{గడుపు}}
\textbf{\te{గడుపు}} — \emph{gaḍupu} — \textquotedblleft{}to spend (time)\textquotedblright{}. Said of time: \textbf{\te{నేను} \te{సెలవులను} \te{గ్రామంలో} \te{గడిపాను}} — \emph{nēnu selavulanu grāmaṁlō gaḍipānu} — \textquotedblleft{}I spent the holidays in the village.\textquotedblright{}

\begin{grammarlens}[title={What you've built}]
One more word for this chapter.
\end{grammarlens}

\subsection*{Your turn}
\begin{itemize}
\raggedright
  \item \emph{Say it:} \emph{gaḍupu}
  \item \emph{Say it:} \emph{gaḍupu}, once more
  \item \emph{Say it:} say \emph{gaḍupu}
  \item \emph{Recall:} say the Telugu for last, then the Telugu for once, then say \emph{gaḍupu} again
\end{itemize}

\begin{tcolorbox}[breakable,colback=teal!4,colframe=teal!35!black,title={Before you move on}]
What does \te{గడుపు} mean? (\textquotedblleft{}To spend (time)\textquotedblright{}.) Say it once more.
\end{tcolorbox}
